% GENERATED FILE. Edit canonical lessons, then run npm run generate:books.
% canonical-source-hash: fnv1a64:3a80dec9cbf584c1
% canonical-lessons: TA-C40-this, TA-C40-that, TA-C40-here, TA-C40-there, TA-C40-who, TA-C40-where, TA-C40-deixis-system

\chapter{Pointing, and Asking}
\label{ch:ta-pointing-and-asking}
\hlchaptermodality{40}

\begin{chapteropening}
\textbf{By the end of this chapter:} \emph{I can point at something in Tamil and ask who or where it is — and I can say what the six words have in common.}

Say the near word, the far word and the question word in a row, and name the part that changed.
\end{chapteropening}

\section[\texorpdfstring{\ta{இது} (idu)}{இது (idu)}]{\ta{இது} (idu) — this one — the thing near me}

\label{lesson:TA-C40-this}



\begin{quote}
You can already name things. This is how you POINT at them.
\end{quote}

\subsection*{You'll want to know: \ta{இது}}
\textbf{\ta{இது}} — \emph{idu} — this one — the thing near me.

Say it, and picture the thing it lands on as you say it. That is the whole word: it does not mean anything on its own, it means whatever the speaker is pointing at.

Its partner, the far one, comes in the next lesson. The two of them differ by one sound at the front — \textbf{i-} for near — and that is not a coincidence. It is the whole system, and the last lesson of this chapter says so.

\textbf{In the family, and next door}

\noindent
\begin{tabularx}{\linewidth}{@{}>{\raggedright\arraybackslash}X>{\raggedright\arraybackslash}X>{\raggedright\arraybackslash}X@{}}
\toprule
\textbf{} & \textbf{word} & \textbf{same root} \\
\midrule
Kannada & \ka{ಇದು} (\emph{idu}) & yes \\
Telugu & \te{ఇది} (\emph{idi}) & yes \\
Malayalam & \ml{ഇത്} (\emph{itŭ}) & yes \\
Hindi (neighbour) & \dv{यह} (\emph{yah}) &  \\
English & this one &  \\
\bottomrule
\end{tabularx}

\subsection*{Your turn}
\begin{itemize}
\raggedright
  \item \emph{Say it:} \textquotedblleft{}\ta{இது}\textquotedblright{} three times, picturing something different each time
  \item \emph{Say it:} it once with a word you already know after it
  \item \emph{Recall:} say \emph{uḍambu}
\end{itemize}

\begin{tcolorbox}[breakable,colback=teal!4,colframe=teal!35!black,title={Before you move on}]
What is \emph{idu}? (this one — the thing near me.) What sound does it start with, and does that sound mean near, far, or a question? (\textbf{i-} — near.)
\end{tcolorbox}
\section[\texorpdfstring{\ta{அது} (adu)}{அது (adu)}]{\ta{அது} (adu) — that one — the thing over there}

\label{lesson:TA-C40-that}



\begin{quote}
You can already name things. This is how you POINT at them.
\end{quote}

\subsection*{You'll want to know: \ta{அது}}
\textbf{\ta{அது}} — \emph{adu} — that one — the thing over there.

Say it, and picture the thing it lands on as you say it. That is the whole word: it does not mean anything on its own, it means whatever the speaker is pointing at.

Its partner is \textbf{\ta{இது}} \emph{idu}, which you have just met. The two of them differ by one sound at the front — \textbf{a-} for far — and that is not a coincidence. It is the whole system, and the last lesson of this chapter says so.

\textbf{In the family, and next door}

\noindent
\begin{tabularx}{\linewidth}{@{}>{\raggedright\arraybackslash}X>{\raggedright\arraybackslash}X>{\raggedright\arraybackslash}X@{}}
\toprule
\textbf{} & \textbf{word} & \textbf{same root} \\
\midrule
Kannada & \ka{ಅದು} (\emph{adu}) & yes \\
Telugu & \te{అది} (\emph{adi}) & yes \\
Malayalam & \ml{അത്} (\emph{atŭ}) & yes \\
Hindi (neighbour) & \dv{वह} (\emph{vah}) &  \\
English & that one &  \\
\bottomrule
\end{tabularx}

\subsection*{Your turn}
Say these aloud:

\begin{itemize}
\raggedright
  \item \textquotedblleft{}\ta{அது}\textquotedblright{} three times, picturing something different each time
  \item it once with a word you already know after it
\end{itemize}

\begin{tcolorbox}[breakable,colback=teal!4,colframe=teal!35!black,title={Before you move on}]
What is \emph{adu}? (that one — the thing over there.) What sound does it start with, and does that sound mean near, far, or a question? (\textbf{a-} — far.)
\end{tcolorbox}
\section[\texorpdfstring{\ta{இங்கே} (iṅgē)}{இங்கே (iṅgē)}]{\ta{இங்கே} (iṅgē) — here — where I am}

\label{lesson:TA-C40-here}



\begin{quote}
You can already name things. This is how you POINT at them.
\end{quote}

\subsection*{You'll want to know: \ta{இங்கே}}
\textbf{\ta{இங்கே}} — \emph{iṅgē} — here — where I am.

Say it, and picture the thing it lands on as you say it. That is the whole word: it does not mean anything on its own, it means whatever the speaker is pointing at.

Its partner, the far one, comes in the next lesson. The two of them differ by one sound at the front — \textbf{i-} for near — and that is not a coincidence. It is the whole system, and the last lesson of this chapter says so.

\textbf{In the family, and next door}

\noindent
\begin{tabularx}{\linewidth}{@{}>{\raggedright\arraybackslash}X>{\raggedright\arraybackslash}X@{}}
\toprule
\textbf{} & \textbf{word} \\
\midrule
Kannada & \ka{ಇಲ್ಲಿ} (\emph{illi}) \\
Telugu & \te{ఇక్కడ} (\emph{ikkaḍa}) \\
Malayalam & \ml{ഇവിടെ} (\emph{iviṭe}) \\
Hindi (neighbour) & \dv{यहाँ} (\emph{yahā̃}) \\
English & here \\
\bottomrule
\end{tabularx}

\subsection*{Your turn}
\begin{itemize}
\raggedright
  \item \emph{Say it:} \textquotedblleft{}\ta{இங்கே}\textquotedblright{} three times, picturing something different each time
  \item \emph{Say it:} it once with a word you already know after it
  \item \emph{Recall:} say \emph{sugam}
  \item \emph{Read it:} \textbf{\ta{எவ்வளவு}}
\end{itemize}

\begin{tcolorbox}[breakable,colback=teal!4,colframe=teal!35!black,title={Before you move on}]
What is \emph{iṅgē}? (here — where I am.) What sound does it start with, and does that sound mean near, far, or a question? (\textbf{i-} — near.)
\end{tcolorbox}
\section[\texorpdfstring{\ta{அங்கே} (aṅgē)}{அங்கே (aṅgē)}]{\ta{அங்கே} (aṅgē) — there — where I am not}

\label{lesson:TA-C40-there}



\begin{quote}
You can already name things. This is how you POINT at them.
\end{quote}

\subsection*{You'll want to know: \ta{அங்கே}}
\textbf{\ta{அங்கே}} — \emph{aṅgē} — there — where I am not.

Say it, and picture the thing it lands on as you say it. That is the whole word: it does not mean anything on its own, it means whatever the speaker is pointing at.

Its partner is \textbf{\ta{இங்கே}} \emph{iṅgē}, which you have just met. The two of them differ by one sound at the front — \textbf{a-} for far — and that is not a coincidence. It is the whole system, and the last lesson of this chapter says so.

\textbf{In the family, and next door}

\noindent
\begin{tabularx}{\linewidth}{@{}>{\raggedright\arraybackslash}X>{\raggedright\arraybackslash}X@{}}
\toprule
\textbf{} & \textbf{word} \\
\midrule
Kannada & \ka{ಅಲ್ಲಿ} (\emph{alli}) \\
Telugu & \te{అక్కడ} (\emph{akkaḍa}) \\
Malayalam & \ml{അവിടെ} (\emph{aviṭe}) \\
Hindi (neighbour) & \dv{वहाँ} (\emph{vahā̃}) \\
English & there \\
\bottomrule
\end{tabularx}

\subsection*{Your turn}
Say these aloud:

\begin{itemize}
\raggedright
  \item \textquotedblleft{}\ta{அங்கே}\textquotedblright{} three times, picturing something different each time
  \item it once with a word you already know after it
\end{itemize}

\begin{tcolorbox}[breakable,colback=teal!4,colframe=teal!35!black,title={Before you move on}]
What is \emph{aṅgē}? (there — where I am not.) What sound does it start with, and does that sound mean near, far, or a question? (\textbf{a-} — far.)
\end{tcolorbox}
\section[\texorpdfstring{\ta{யார்} (yār)}{யார் (yār)}]{\ta{யார்} (yār) — who? — asking about a person}

\label{lesson:TA-C40-who}



\begin{quote}
You can already name things. This is how you ASK about them.
\end{quote}

\subsection*{You'll want to know: \ta{யார்}}
\textbf{\ta{யார்}} — \emph{yār} — who? — asking about a person.

Say it with an asking voice. It points at no one: it asks your listener to name the person.

\emph{yār} opens with \textbf{yā-}, not with the \textbf{i-} of near or the \textbf{a-} of far that the four words before it share. It is the one word in this chapter outside that pattern, so learn it as itself: the asking word for people.

\textbf{In the family, and next door}

\noindent
\begin{tabularx}{\linewidth}{@{}>{\raggedright\arraybackslash}X>{\raggedright\arraybackslash}X>{\raggedright\arraybackslash}X@{}}
\toprule
\textbf{} & \textbf{word} & \textbf{same root} \\
\midrule
Kannada & \ka{ಯಾರು} (\emph{yāru}) & yes \\
Telugu & \te{ఎవరు} (\emph{evaru}) & yes \\
Malayalam & \ml{ആര്} (\emph{ārŭ}) & yes \\
Hindi (neighbour) & \dv{कौन} (\emph{kaun}) &  \\
English & who? &  \\
\bottomrule
\end{tabularx}

\subsection*{Your turn}
\begin{itemize}
\raggedright
  \item \emph{Say it:} \textquotedblleft{}\ta{யார்}\textquotedblright{} three times, asking about a different person each time
  \item \emph{Say it:} it once with a word you already know after it
  \item \emph{Recall:} say \emph{viḍai}
  \item \emph{Read it:} \textbf{\ta{ஒரு}}
\end{itemize}

\begin{tcolorbox}[breakable,colback=teal!4,colframe=teal!35!black,title={Before you move on}]
What is \emph{yār}? (who? — asking about a person.) What sound does it start with: the i- of near, the a- of far, or neither? (\textbf{Neither} — \textbf{yā-}; it asks about a person.)
\end{tcolorbox}
\section[\texorpdfstring{\ta{எங்கே} (eṅgē)}{எங்கே (eṅgē)}]{\ta{எங்கே} (eṅgē) — where? — asking about a place}

\label{lesson:TA-C40-where}



\begin{quote}
You can already name things. This is how you ASK about them.
\end{quote}

\subsection*{You'll want to know: \ta{எங்கே}}
\textbf{\ta{எங்கே}} — \emph{eṅgē} — where? — asking about a place.

Say it with an asking voice. It points at no place: it asks your listener to point one out.

Its partners are \textbf{\ta{இங்கே}} \emph{iṅgē} and \textbf{\ta{அங்கே}} \emph{aṅgē}, which you have already met. The three differ by one sound at the front — \textbf{e-} for ask — and that is not a coincidence. It is the whole system, and the last lesson of this chapter says so.

\textbf{In the family, and next door}

\noindent
\begin{tabularx}{\linewidth}{@{}>{\raggedright\arraybackslash}X>{\raggedright\arraybackslash}X@{}}
\toprule
\textbf{} & \textbf{word} \\
\midrule
Kannada & \ka{ಎಲ್ಲಿ} (\emph{elli}) \\
Telugu & \te{ఎక్కడ} (\emph{ekkaḍa}) \\
Malayalam & \ml{എവിടെ} (\emph{eviṭe}) \\
Hindi (neighbour) & \dv{कहाँ} (\emph{kahā̃}) \\
English & where? \\
\bottomrule
\end{tabularx}

\subsection*{Your turn}
Say these aloud:

\begin{itemize}
\raggedright
  \item \textquotedblleft{}\ta{எங்கே}\textquotedblright{} three times, asking about a different place each time
  \item it once with a word you already know after it
\end{itemize}

\begin{tcolorbox}[breakable,colback=teal!4,colframe=teal!35!black,title={Before you move on}]
What is \emph{eṅgē}? (where? — asking about a place.) What sound does it start with, and does that sound mean near, far, or a question? (\textbf{e-} — ask.)
\end{tcolorbox}
\section[i- / a- / e-]{i- / a- / e- — six words, one machine}

\label{lesson:TA-C40-deixis-system}



\begin{quote}
Say the six words from this chapter out loud, in order. Listen to how they begin.
\end{quote}

\begin{grammarlens}[title={Grammar lens: one slot changes}]
Five of these six words are \textbf{three beginnings}, reused; only the person question stands apart.

\noindent
\begin{tabularx}{\linewidth}{@{}>{\raggedright\arraybackslash}X>{\raggedright\arraybackslash}X@{}}
\toprule
\textbf{meaning} & \textbf{word} \\
\midrule
near & \textbf{\ta{இங்கே}} \emph{iṅgē} \\
far & \textbf{\ta{அங்கே}} \emph{aṅgē} \\
question & \textbf{\ta{எங்கே}} \emph{eṅgē} \\
\bottomrule
\end{tabularx}

Tamil shows the system at its cleanest: \ta{இங்கே}, \ta{அங்கே}, \ta{எங்கே} differ in one letter.

Change the front, and the meaning walks from \textbf{near} to \textbf{far} to \textbf{a question}. Nothing else in the word moves.

The same three beginnings run the thing-words too: \textbf{\ta{இது}} \emph{idu} and \textbf{\ta{அது}} \emph{adu} start exactly like the first two rows. This chapter did not teach you the matching question word for things — that one arrives later, and when it does you will already know what it has to start with.

This is worth more than the six words are. When you meet a new word in this family later, you will not have to be taught all three of it — you will be taught one and work out the others.
\end{grammarlens}

\subsection*{Your turn}
\begin{itemize}
\raggedright
  \item \emph{Say it:} the near word, then the far word, then the question word
  \item \emph{Say it:} them again for place, not for things
  \item \emph{Notice:} the part that did NOT change
\end{itemize}

\begin{tcolorbox}[breakable,colback=teal!4,colframe=teal!35!black,title={Before you move on}]
Which beginning means near? (\textbf{i-}) Far? (\textbf{a-}) And which one turns the word into a question? (\textbf{e-})
\end{tcolorbox}
